\documentclass[11pt]{amsart}
\usepackage[T1]{fontenc}
\usepackage{lmodern,microtype,amsmath,amssymb,amsthm,mathtools,booktabs}
\usepackage[margin=1.1in]{geometry}
\usepackage[hidelinks]{hyperref}
\usepackage{xurl}

\newtheorem{theorem}{Theorem}[section]
\newtheorem{lemma}[theorem]{Lemma}
\newtheorem{corollary}[theorem]{Corollary}
\theoremstyle{remark}
\newtheorem{remark}[theorem]{Remark}
\numberwithin{equation}{section}

\newcommand{\E}{\mathbb{E}}
\renewcommand{\P}{\mathbb{P}}

\hypersetup{pdftitle={Occupation sums of the harmonic descent chain and slowly decaying weights},pdfauthor={Yitzchak Shmalo}}

\title[Occupation sums and slowly decaying weights]{Occupation sums of the harmonic descent chain and slowly decaying weights}
\author{Yitzchak Shmalo}
\address{Einstein Institute of Mathematics, The Hebrew University of Jerusalem, Jerusalem, Israel}
\email{yitzchak.shmalo@gmail.com}
\date{7 November 2026}

\begin{document}

\begin{abstract}
Let $a(n,i)$ be the probability that the harmonic descent chain started at $n$ visits $i$, and $a(i)=\lim_n a(n,i)$.
Aldous and Janson asked whether $\sum_{i=2}^n a(n,i)f(i)\to\sum_{i\ge2}a(i)f(i)$ for every non-negative $f$ of
polynomial growth with a finite limit sum, and Brandenberger, Chin and Mossel proved this when $f(i)=O(i^{-\varepsilon})$.
We show that for every $\delta>0$ there is an $f$ with $f(i)=O((\log i)^{-1+\delta})$ and
$\sum_i a(i)f(i)<\infty$ for which the occupation sums are unbounded, along values of $n$ with $f(n)=0$. The cause is that
the chain started at $n$ makes, on average, at least of order $\log n/\log\log n$ visits to a window of width
$n/(\log n)^2$ just below $n$.
\end{abstract}

\maketitle

\section{The question}

The harmonic descent chain is the Markov chain on $\{1,2,\dots\}$ which, from a state $m\ge2$, jumps to $m-d$ with
probability
\begin{equation}\label{eq:kernel}
p(m,m-d)=\frac{1}{d\,h_{m-1}},\qquad 1\le d\le m-1,\qquad h_r=\sum_{j=1}^r\frac1j,
\end{equation}
and stays at $1$ forever. Started at $n$, it records the clade sizes on the path from the root to a uniformly chosen leaf
of the discrete-time critical beta-splitting tree with $n$ leaves \cite{AP,AJ,AJL}. Let $a(n,i)$ be the probability that
the chain started at $n$ ever visits $i$, so that $a(n,n)=a(n,1)=1$. For fixed $i\ge2$ the limit
\begin{equation}\label{eq:limit}
a(i)=\lim_{n\to\infty}a(n,i)=\frac{6\,h_{i-1}}{\pi^2(i-1)}
\end{equation}
exists \cite[Theorem~18]{AJ}, \cite{AJL}. Aldous and Janson \cite[Open Problem~8]{AJ} asked for bounds on
$|a(n,i)-a(i)|$, and in particular for a proof of the following statement.

\medskip\noindent\textbf{Ansatz 20 of \cite{AJ}.} For a non-negative sequence $(f(j),\,j\ge2)$ with $f(j)=O(j^k)$ for some
$k<\infty$:
\begin{enumerate}
\item[(i)] if $\sum_{i\ge2}a(i)f(i)<\infty$ then $\sum_{i=2}^n a(n,i)f(i)\to\sum_{i\ge2}a(i)f(i)$;
\item[(ii)] if $\sum_{i\ge2}a(i)f(i)=\infty$ then $\sum_{i=2}^n a(n,i)f(i)\sim\sum_{i=2}^n a(i)f(i)$.
\end{enumerate}

\medskip
Brandenberger, Chin and Mossel \cite{BCM} proved (i) for every non-negative $f$ with $f(i)=O(i^{-\varepsilon})$ for some
$\varepsilon>0$, as a consequence of their rate $a(n,i)-a(i)=n^{-\beta+o(1)}$; Svihla \cite{S} gives the exact
asymptotic expansion of $a(n,i)-a(i)$ for $i=o(\sqrt n)$. Since $a(n,n)=1$, part (i) as stated fails for any $f$ with
$\sum a(i)f(i)<\infty$ and $\limsup f(n)>0$, through the single term $i=n$; Aldous and Janson remark that the method is
meant for statistics that are not sensitive to the near-root structure. The point of this note is that the near-root
structure matters well beyond that one term: the states just below $n$ are visited many times, and a weight that decays
only like a power of $\log i$ can see them.

\begin{theorem}\label{thm:main}
For every $\delta\in(0,1)$ there is a sequence $f\ge0$ with $f(i)=O((\log i)^{-1+\delta})$ and
$\sum_{i\ge2}a(i)f(i)<\infty$, and integers $N_j\to\infty$ with $f(N_j)=0$, such that
\[
\sum_{i=2}^{N_j}a(N_j,i)f(i)\longrightarrow\infty .
\]
In particular part (i) of Ansatz 20 fails for this $f$, and it fails for the sum over $2\le i\le n-1$ as well.
\end{theorem}

The decay $f(i)=O((\log i)^{-1+\delta})$ cannot be replaced by $O(i^{-\varepsilon})$, by \cite{BCM}. Between the two
the question is open; see Remark~\ref{rem:gap}.

\section{Visits just below the starting point}

The chain is strictly decreasing until it reaches $1$, so it visits each state at most once and, for a set $B$ of states
in $\{2,\dots,n-1\}$,
\begin{equation}\label{eq:visits}
\sum_{i\in B}a(n,i)=\E_n\bigl[\#\{t\ge1:\ X_t\in B\}\bigr].
\end{equation}

\begin{lemma}\label{lem:band}
Let $n\ge4$ and let $T,L$ be integers with $1\le T\le L\le n-2$. Put $x=\lfloor L/T\rfloor$ and $q=h_x/h_{n-1}$. Then
\[
\sum_{i=n-L}^{n-1}a(n,i)\ \ge\ \sum_{t=1}^{T}q^t\ \ge\ T\,q^T .
\]
\end{lemma}

\begin{proof}
From any state $m$ with $2\le m\le n$, by \eqref{eq:kernel} the jump $d$ satisfies
$\P(d\le x)=h_{\min(x,m-1)}/h_{m-1}\ge h_x/h_{n-1}=q$: if $x\le m-1$ this holds because $h_{m-1}\le h_{n-1}$, and if
$x\ge m-1$ the probability is $1$. Let $E_t$ be the event that each of the first $t$ jumps is at most $x$. On $E_t$, for
$1\le s\le t\le T$, we have $n-1\ge X_s\ge n-sx\ge n-Tx\ge n-L\ge2$, so the chain has not been absorbed and $X_s$ lies in
$B=\{n-L,\dots,n-1\}$. By the Markov property $\P_n(E_t)\ge q^t$. Hence the number of visits to $B$ is at least
$\sum_{t=1}^T\mathbf 1_{E_t}$, and \eqref{eq:visits} gives the first inequality. The second holds since $q\le1$.
\end{proof}

\begin{corollary}\label{cor:many}
Let $N_j=2^{j^2}$, $L_j=\lfloor N_j/j^4\rfloor$ and $B_j=\{N_j-L_j,\dots,N_j-1\}$. Then for all large $j$,
\[
\sum_{i\in B_j}a(N_j,i)\ \ge\ e^{-9}\Bigl\lfloor\frac{j^2}{\log j}\Bigr\rfloor .
\]
\end{corollary}

\begin{proof}
Write $\ell=\log N_j=j^2\log2$ and take $T=\lfloor j^2/\log j\rfloor$, so $T\le L_j$ for large $j$. With
$x=\lfloor L_j/T\rfloor$ we have $x+1\ge L_j/T\ge N_j/(2j^4T)$, so, using $h_x\ge\log(x+1)$, $T\le j^2$ and
$h_{N_j-1}\le1+\ell$,
\[
q=\frac{h_x}{h_{N_j-1}}\ \ge\ \frac{\ell-4\log j-\log T-\log 2}{1+\ell}\ \ge\ 1-u,\qquad u=\frac{2+6\log j}{1+\ell}.
\]
Then $u\to0$ and $Tu\le(2+6\log j)/(\log j\cdot\log2)\to6/\log2<8.7$, so
$q^T\ge\exp\bigl(-Tu/(1-u)\bigr)\ge e^{-9}$ for large $j$, and Lemma~\ref{lem:band} gives the bound.
\end{proof}

\section{Proof of Theorem~\ref{thm:main}}

Fix $\delta\in(0,1)$ and take $N_j$, $L_j$, $B_j$ as in Corollary~\ref{cor:many}. Choose $j_0$ so large that the
corollary holds for $j\ge j_0$ and $N_j-L_j>N_{j-1}$ for $j>j_0$; the blocks $B_j$, $j\ge j_0$, are then disjoint. Set
\[
f(i)=j^{-2+2\delta}\quad(i\in B_j,\ j\ge j_0),\qquad f(i)=0\ \text{otherwise}.
\]
Then $f(N_j)=0$ for every $j$.

\emph{Decay.} For $i\in B_j$ we have $\log i\le\log N_j=j^2\log2$, so
\[
f(i)=(j^2)^{-(1-\delta)}=(\log 2)^{1-\delta}(\log N_j)^{-(1-\delta)}\le(\log i)^{-(1-\delta)}.
\]

\emph{Finite limit sum.} By \eqref{eq:limit} and $h_{i-1}\le1+\log i$ there is a constant $C$ with
$a(i)\le C\log i/i$ for $i\ge2$. For $i\in B_j$ we have $i\ge N_j/2$, so
\[
\sum_{i\in B_j}a(i)\ \le\ C\,(j^2\log2)\,\frac{2\,(L_j+1)}{N_j}\ \le\ \frac{3C}{j^{2}}
\]
for large $j$, and $\sum_i a(i)f(i)\le\sum_j 3C\,j^{-4+2\delta}<\infty$.

\emph{Unbounded occupation sums.} By Corollary~\ref{cor:many},
\[
\sum_{i=2}^{N_j}a(N_j,i)f(i)\ \ge\ j^{-2+2\delta}\sum_{i\in B_j}a(N_j,i)\ \ge\ e^{-9}\,j^{-2+2\delta}
\Bigl\lfloor\frac{j^2}{\log j}\Bigr\rfloor\ \longrightarrow\ \infty .
\]
Every term used has $i\le N_j-1$, which proves the last sentence of the theorem.\qed

\begin{remark}
The weight in the proof is bounded, so it satisfies the growth condition of Ansatz 20 with $k=0$. Taking
$f=\mathbf 1_{\bigcup_j B_j}$ (the case $\delta=1$) gives $\sum_i a(i)f(i)<\infty$ by the same estimate and occupation
sums of order at least $j^2/\log j$ at $N_j$.
\end{remark}

\begin{remark}\label{rem:gap}
The proof uses a window of relative width $(\log n)^{-2}$. A block of weight $c$ on a window of relative width $w$
contributes about $c\,w\log n$ to the limit sum, while Lemma~\ref{lem:band} gives it at least $c\log n/\log\log n$ in the
occupation sum when $w$ is at least a power of $(\log n)^{-1}$; for a window $n^{1-\eta}$ the jumps out of the window are
no longer rare and the lemma gives only a bounded number of visits. The construction therefore stops short of weights of
order $(\log i)^{-1}$, and we do not know whether part (i) holds for every $f$ with $f(i)=O((\log i)^{-1})$, or for
every $f$ with $f(i)\to0$ that is monotone or regularly varying. For monotone $f$ the window just below $n$ carries weight close to $f(n)$, and the question becomes one
about the sizes of $a(n,i)$ for $i$ near $n$, a range not covered by \cite{BCM} or \cite{S}.
\end{remark}

\section{A computation}

By the Markov property, $a(n,n)=1$ and $a(n,i)=\sum_{m=i+1}^{n}a(n,m)\,/\,\bigl((m-i)\,h_{m-1}\bigr)$ for $i<n$; for $i$ in
a window $\{n-L,\dots,n-1\}$ only states in the window and $n$ itself contribute, so the window sums can be computed
exactly in $O(L^2)$ operations for very large $n$. Table~\ref{tab:band} lists $\sum_{i\in B_j}a(N_j,i)$ for the blocks of
Corollary~\ref{cor:many} and the sums over the top sixteenth $\{n-n/16,\dots,n-1\}$ for $n=2^{10},2^{14},2^{18}$. The
latter grow by about $0.88$ each time $n$ is multiplied by $16$, in line with the count of order $\log n$ in
Remark~\ref{rem:gap}.

\begin{table}[ht]
\centering
\begin{tabular}{rrrr}
\toprule
$n$ & window $L$ & $\sum_{i=n-L}^{n-1}a(n,i)$ & \\
\midrule
$2^{9}$ & $\lfloor n/3^4\rfloor=6$ & $0.479$ & $j=3$\\
$2^{16}$ & $n/4^4=256$ & $1.022$ & $j=4$\\
$2^{25}$ & $\lfloor n/5^4\rfloor=53{,}687$ & $1.696$ & $j=5$\\
\midrule
$2^{10}$ & $n/16$ & $1.407$ & \\
$2^{14}$ & $n/16$ & $2.286$ & \\
$2^{18}$ & $n/16$ & $3.172$ & \\
\bottomrule
\end{tabular}
\caption{Expected visits of the harmonic descent chain started at $n$ to the window just below $n$.}\label{tab:band}
\end{table}

\section*{Acknowledgements and funding}
The research presented in this paper was supported by the European Research
Council (ERC) under the European Union's Horizon Europe research and
innovation programme (grant agreement No.~101041711), by the Simons
Foundation as part of the Collaboration on the Mathematical and Scientific
Foundations of Deep Learning, by Heights Labs, by Convex Nexus Capital, by
the Israel Science Foundation (grant number 2258/19), and by the Israel
Science Foundation (ISF Grant 4101/25). Computational resources were
provided in part by Google Cloud through its education credits program
(credit 530550103).

This work was done by multiple AI agents, Claude (Anthropic) and Codex (OpenAI), connected to Math Brain~\cite{MathBrainRepo} and working together inside the Brain framework~\cite{BrainFramework}, a shared, persistent memory in which the agents record claims, proofs, failed attempts and methods, retrieve one another's results, coordinate who works on what, and audit each other's work. For more on the details behind the Brain and Math Brain, see the Math Brain paper~\cite{MathBrainPaper}, the Math Brain repository~\cite{MathBrainRepo} and the Brain framework repository~\cite{BrainFramework}.

\section*{Declaration of competing interest}
The author holds equity in Heights Labs and in Convex Nexus Capital and
has a management role at Convex Nexus Capital. These entities are included
in the funding declaration above.

\section*{Declaration of generative AI and AI-assisted technologies}
The paper comes out of Math Brain, an ongoing project on the use of AI systems for research mathematics. In November 2026
Claude Opus~5.5 (Anthropic), working in Claude Code, found the construction and the proofs, wrote and ran the
computation of Table~\ref{tab:band}, and wrote the paper. The author directed the work and is responsible for the
content of the paper and for decisions concerning its dissemination.

\section*{Data availability}
No data were used. The code for Table~\ref{tab:band}, with its output, is available at
\url{https://github.com/yspennstate/harmonic-descent-occupation-sums}.

The paper is part of the Math Brain project, whose public repository is \url{https://github.com/yspennstate/math-brain-public}~\cite{MathBrainRepo}.

\begin{thebibliography}{9}
\bibitem{AJ} D.~Aldous and S.~Janson, \emph{The critical beta-splitting random tree II: overview and open problems},
arXiv:2303.02529v3, 2024.
\bibitem{AJL} D.~Aldous, S.~Janson and X.~Li, \emph{The harmonic descent chain}, arXiv:2405.05102, 2024.
\bibitem{AP} D.~Aldous and B.~Pittel, \emph{The critical beta-splitting random tree I: heights and related results},
arXiv:2302.05066.
\bibitem{BCM} A.~Brandenberger, B.~Chin and E.~Mossel, \emph{Asymptotics for the harmonic descent chain and applications
to critical beta-splitting trees}, arXiv:2505.24821v3, 2026.
\bibitem{S} S.~P.~Svihla, \emph{Exact asymptotics for the harmonic descent chain and extended convergence bounds for
recursive statistics}, arXiv:2611.00347, 2026.
\bibitem{MathBrainRepo} Y.~Shmalo, \emph{Math Brain: public snapshot, code and tests}, 2026, \url{https://github.com/yspennstate/math-brain-public}.

\bibitem{MathBrainPaper} Y.~Shmalo, \emph{Math Brain: shared, auditable memory for multi-agent mathematical research}, October 2026, \url{https://github.com/yspennstate/math-brain-public/tree/main/papers/math-brain-2026-10-24}.

\bibitem{BrainFramework} Y.~Shmalo, \emph{Brain Framework}, public software repository, 2026, \url{https://github.com/yspennstate/brain-framework-public}.

\end{thebibliography}

\end{document}
